\scanpage{59}
% archive foliation: 61
\uline{Th. 4} Équivalence de catégories entre

a) Cat. des couples $(X,K,S)$ $X$ surface conforme
à bord compacte, $(K,S)$ complexe top. \deleted{\unsure{cpt}}\inserted{cpt} avec
\[
\delta X\subset K\subset X,\qquad K\text{ assez régulier}\quad\text{(rectifiable suffit)}
\]

\margin{\uline{NB} On ne suppose pas les \unsure{compos. conn.}
de $K\setminus S$ des intervalles
(il peut y avoir des cercles).}

b) Cat des systèmes $(\widetilde X,\widetilde S,\sigma)$, où
$\widetilde X$ est une surface conforme à bord \inserted{compacte},
$\widetilde S$ une partie finie de $\delta\widetilde X$, $\sigma$ une
involution de $\operatorname{\text{Déc}}(\delta\widetilde X/\widetilde S)$, $\sigma$ assez
\unsure{régulière|régulier} (rectifiable suffit).

Le foncteur a) $\longrightarrow$ b) résulte du th. $\unsure{3|2}$,
avec un complément qui explicite les
propriétés de régularité de la fonction-recollement
dont il a été question dans
le lemme 2. Pour le foncteur en sens
inverse, il est connu topologiquement, et
comme faisceau structural conforme
(large) sur $X$ \inserted{$=\widetilde X/R$}, on prend les fonctions
continues qui sont conformes dans $X\setminus K\simeq\operatorname{Int}\widetilde X$.
